\chapter{\ifenglish Conclusion\else บทสรุป\fi}

\section{สรุปและบทสรุป}


โครงการ eDENT-ELDER Phase 2 มีเป้าหมายเพื่อแก้ไขข้อจำกัดของระบบเดิม
และพัฒนาแพลตฟอร์มสำหรับการคัดกรองมะเร็งช่องปากในระดับชุมชนให้สามารถใช้งานได้จริง

จากผลการดำเนินงานในบทก่อนหน้า พบว่าระบบสามารถรองรับกระบวนการคัดกรองแบบครบวงจร
ตั้งแต่การเก็บข้อมูลผู้ป่วย การวิเคราะห์ด้วย AI ไปจนถึงการให้คำปรึกษาจากผู้เชี่ยวชาญ
โดยมีความเสถียรและประสิทธิภาพในระดับที่เหมาะสมต่อการใช้งานจริง

นอกจากนี้ การออกแบบระบบในลักษณะ Modular Monolith architecture
ร่วมกับแนวคิด Dual-Path AI และ DevOps Pipeline
ช่วยให้ระบบมีความยืดหยุ่น สามารถขยายระบบได้ในอนาคต
และลดความเสี่ยงจากความล้มเหลวของระบบ

\begin{tcolorbox}[bluebox, title=ผลสำเร็จหลัก]
  \begin{itemize}[leftmargin=1.4em, itemsep=5pt]
    \item \textbf{สถาปัตยกรรม:}
          Modular Monolith architecture (Next.js + FastAPI + DDD) แทนสถาปัตยกรรมเดิมแบบ monolith
    \item \textbf{ระบบ AI:}
          Dual-Path AI มี inference latency $<200\,\text{ms}$ ภายใต้ synthetic load
          อัตรา request สำเร็จ 99.46\% จากทั้งหมด 7,632 requests
    \item \textbf{ความปลอดภัย:}
          กรอบ PDPA/HIPAA พร้อม Tamper-Proof Audit Logging
          และ Role-Based Access Control (RBAC) ครอบคลุม 6 บทบาท
    \item \textbf{DevOps:}
          Docker Containerization ครบวงจร พร้อม CI/CD Pipeline
          และการแยก Environment เข้มงวด
    \item \textbf{คุณภาพ:}
          900 automated tests ผ่าน 100\%
          100\% coverage บน critical business logic
    \item \textbf{ความร่วมมือ:}
          หลายสถาบัน ครอบคลุมทันตแพทยศาสตร์ วิศวกรรมศาสตร์ และสาธารณสุขชุมชน
  \end{itemize}
\end{tcolorbox}

แพลตฟอร์มพร้อม Deploy ระดับนำร่องในเชียงใหม่ พะเยา ภูเก็ต และหนองบัวลำภู
เพื่อช่วยให้ตรวจพบมะเร็งช่องปากในผู้สูงอายุชนบทเร็วขึ้น
ลดการวินิจฉัยระยะท้าย และพัฒนาผลลัพธ์ผู้ป่วยในวงกว้าง

โครงการนี้ประสบความสำเร็จในการแก้ปัญหาคอขวดและจุดล้มเหลวของระบบเดิม
เปลี่ยนผ่านสู่ \textbf{แพลตฟอร์มทางการแพทย์พร้อมใช้งานจริง (Production-Ready)}
ด้วยสถาปัตยกรรมที่เสถียร ปลอดภัย และเหมาะสมต่อการใช้งานจริงในระดับบริการ
\vspace{0.5em}
\begin{tcolorbox}[graybox]
  \textbf{ข้อจำกัดของโครงการ (Limitations):}
  \begin{itemize}[leftmargin=1.4em, itemsep=3pt]
    \item ระบบพึ่งพา AI Model จากภายนอก
          ซึ่งอาจมีข้อจำกัดด้านการควบคุมคุณภาพและการปรับปรุงโมเดล
    \item การทดสอบยังอยู่ในระดับนำร่อง
          และยังไม่ได้ประเมินผลในระดับประชากรขนาดใหญ่
    \item ประสิทธิภาพของระบบขึ้นอยู่กับคุณภาพของภาพที่ผู้ใช้งานส่งเข้ามา
    \item การยอมรับของผู้ใช้งานในระยะยาวยังต้องมีการศึกษาเพิ่มเติม
  \end{itemize}
\end{tcolorbox}
\vspace{0.5em}
\begin{tcolorbox}[graybox]
  \textbf{งานในอนาคต (Future Work):}
  \begin{itemize}[leftmargin=1.4em, itemsep=3pt]
    \item ขยาย AI Model ครอบคลุมรอยโรคช่องปากเพิ่มเติม
          และปรับปรุงความแม่นยำด้วย Dataset Annotation จากพื้นที่นำร่อง
    \item ขยายสู่จังหวัดเพิ่มเติมและเชื่อมต่อกับระบบสารสนเทศสุขภาพ (HIS) ของไทย
    \item ดำเนินการศึกษา Clinical Validation อย่างเป็นทางการ
          วัด Stage-Shift และเวลาพบผู้เชี่ยวชาญในกลุ่มนำร่อง
    \item สำรวจการเชื่อมต่อ LINE OA และ PWA Offline Mode
          สำหรับชุมชนที่มีอินเทอร์เน็ตจำกัด
  \end{itemize}
\end{tcolorbox}
\vspace{0.5em}
โดยสรุป โครงการนี้แสดงให้เห็นถึงศักยภาพของการผสานเทคโนโลยีดิจิทัล
เข้ากับระบบสาธารณสุขในระดับชุมชนอย่างเป็นรูปธรรม
และสามารถต่อยอดสู่การใช้งานในระดับประเทศได้ในอนาคต
\vfill
{\color{primaryblue}\rule{\linewidth}{1pt}}
\begin{center}
  \small\color{medgray}
  eDENT-ELDER เฟส 2 รายงานโครงงานฉบับสมบูรณ์ $\cdot$
  รหัสวิชา 261492 $\times$ 2568 $\cdot$
  ภาควิชาวิศวกรรมคอมพิวเตอร์ คณะวิศวกรรมศาสตร์ มหาวิทยาลัยเชียงใหม่
\end{center}
